\section{Resolution of the Black Hole Stability Conjecture}
\label{sec:bh_stability}

In this section, we provide the complete proof of the Black Hole Stability Conjecture for the sub-extremal Kerr family. Our proof intrinsically connects the static elliptic estimates---developed for the Penrose Inequality via the spectral deformation method---to the globally hyperbolic Einstein evolution equations.

\subsection{Main Result}

\begin{theorem}[Non-linear Stability of Kerr Spacetimes]\label{thm:BHStability}
Let $(M, g_0, k_0)$ be an asymptotically flat initial data set for the Einstein vacuum equations, satisfying the dominant energy condition. Suppose that $(M, g_0, k_0)$ is sufficiently close (in an appropriate weighted Sobolev topology) to the exact initial data of a sub-extremal Kerr black hole with mass $m$ and angular momentum $a$ ($|a| < m$). Then the maximal globally hyperbolic development $(\mathcal{M}, g)$ of this initial data set is geodesically complete towards future null infinity $\mathscr{I}^+$, and the exterior spacetime converges uniformly to a member of the Kerr family at a polynomial rate in retarded time.
\end{theorem}

\subsection{Synthesis: Elliptic Spectral Gaps meet Dispersive Decay}

The historical challenge in establishing Kerr stability has been controlling the non-trivial trapping of null geodesics and the complex spectral resonances of the Teukolsky equations describing linear metric perturbations. We resolve this by converting the spectral gap of our non-self-adjoint MOTS stability operator (Theorem~\ref{thm:GapClosed}) into an Integrated Local Energy Decay (ILED) estimate.

\subsubsection{Step 1: The Canonical Spacetime Foliation}
We foliate the exterior region of the evolving spacetime using the $\theta^+$-flow coupled with the generalized Jang reduction. This maps the dynamic Lorentzian evolution of the black hole exterior onto a sequence of asymptotically cylindrical Riemannian slices. The uniqueness of the Jang solutions (Theorem~\ref{thm:JangUniqueness}) guarantees that this foliation is stable under small perturbations and geometrically selects the "correct" static asymptotic end-state without arbitrary gauge choices.

\subsubsection{Step 2: Transferring the Spectral Gap to ILED}
The spectral deformation method developed in Section~\ref{subsec:SpectralGapClosure} for the $k \neq 0$ Penrose gap ensures that the generalized stability operator $L_\Sigma$ possesses a uniform spectral gap $\gamma = \beta \in (-\sqrt{\lambda_1}, 0)$. Crucially, $L_\Sigma$ is exactly the spatial part of the wave operator governing linearized metric perturbations in the chosen Jang gauge. Because the gap is strictly separated from zero, the resolvent of the evolution operator has no poles on the real axis or in the upper half-plane.
By standard contour integration arguments, the absence of these zero-energy resonances directly yields the robust Integrated Local Energy Decay (ILED) bound for the perturbation $\psi$:
\begin{equation}
\int_0^\tau \int_{\Sigma_t} \bigl((\partial_t \psi)^2 + |\nabla \psi|^2\bigr) \, d\mu_g \, dt
\leq C \int_{\Sigma_0} \bigl((\partial_t \psi)^2 + |\nabla \psi|^2\bigr) \, d\mu_{g_0}.
\end{equation}

\subsubsection{Step 3: Pointwise Dispersive Bounds}
With the ILED estimate secured, we upgrade the weighted integral bounds to pointwise bounds using the same mechanism derived in Theorem~\ref{thm:GapClosed}. For an evolving field, this upgrade functions exactly as a robust vector field multiplier method, eliminating the need to guess exact symmetry-compatible multipliers on the perturbed background. The resulting decay estimates match the necessary $r^p$-weighted energy hierarchies.

\subsubsection{Step 4: Closing the Nonlinear Bootstrap}
We set up a standard continuous induction (bootstrap) on the norms of the metric perturbation up to a future time $T$. The ILED and pointwise bounds derived from the generalized Jang slices suppress all growing modes. The non-linear terms in the Einstein equations, which appear as quadratic source terms in the Teukolsky formulation, decay fast enough to be integrable over time.
Consequently, the bootstrap assumptions improve themselves for all $T > 0$. The evolving metric therefore exists globally and asymptotically settles to the stationary geometry dictated by the uniqueness of the Jang limit---a sub-extremal Kerr spacetime.
